\section{Conclusion}
\label{sec:conclusion}

GPU partitioning isolates resources; it does not schedule a dependency.
\sys turns the producer's activation lifecycle into a cross-partition
contract.  Its full-coherent registered system-memory activation edge carries
the same physical payload between separate Thor MIG CUDA contexts, its planner
admits only a hash-bound joint-tail reservation, and its runtime returns
best-effort capacity immediately after publication.  The production wall
includes operational release, queueing, and protection, while correctness
validation remains independently observable after completion.

In six counterbalanced, thermal-normalized ImageNette sessions, \sys completes
6,600 requests with zero misses, a $1{,}902.987~\mu$s p99, and a 0.0454\%
one-sided 95\% DMR bound at the frozen $2{,}255.483~\mu$s deadline.  Relative
to NVIDIA MPS, its paired session-p99 reduction is $138.508~\mu$s and its
background-goodput difference is unresolved.  A separate directional gate
repeats \sys, NVIDIA MIG, NVIDIA MPS, native XSched, Orion, and native Pantheon
on the same 90 labelled inputs, arrival trace, and frozen deadline.  ImageNette
and LibriSpeech ASR preserve their reference metrics.  BLESS, BOER, ParvaGPU,
gpulet, GSLICE, and DeepPlan remain visible in a separate partial-evidence
table, without treating unlike contracts as a common rank.

A separate nonthermal three-slot Whisper motivation isolates where MIG ceases
to be sufficient.  At 19 requests/s, colocating both critical stages on 2g
preserves the isolated 1g tenant but incurs 167/300 misses; static MPS splitting
incurs 64/300; \sys's split and publication-scoped protection incurs 0/300.
The inputs and outputs are real and byte-identical, but the arrivals cyclically
replay twelve windows, so this result establishes a mechanism and load regime
rather than a deployment SLO.

Our claims are bounded to one Thor and a fixed 1g+2g placement.  The formal
vision path is low-inflight, and the three-slot ASR path lacks the standalone
ring's fault-recovery contract; causal, transport, and load controls remain
exploratory or descriptive.  Only \sys is a proposed system name; comparators
retain their original roles.  Within this boundary, the evidence suggests a broader lesson:
on a coherent edge SoC, GPU management should follow the activation, its
publication, and the remaining end-to-end slack connecting dependent stages.
